\section{Introduction}
\label{sec:introduction}

Weak gravitational lensing distorts the apparent shapes of distant galaxies,
producing correlations that trace intervening matter. These correlations
allow astronomers to study
dark matter and test how structure grows as the Universe expands
\citep{2001PhR...340..291B}. Their interpretation, however, requires knowing
which part of the measured signal comes from lensing.

Galaxies are not oriented independently of their surroundings. Their
formation and evolution can produce correlated intrinsic shapes, known as
\emph{intrinsic alignments} (IA). These correlations enter the same
measurements as lensing and can alter the inferred cosmological parameters
if they are modelled inadequately \citep{2025A&ARv..33....5C}. This challenge
becomes more demanding as surveys measure galaxy shapes with greater
statistical precision.

A flexible IA model can accommodate changes with spatial scale, cosmic
time, and galaxy population. Yet additional parameters can be difficult to
constrain, and different combinations may produce similar predictions.
This suggests a specific research question: how much of a flexible
alignment response can be represented with fewer coordinates, and which
features are lost as the representation becomes shorter? A successful
compression would preserve the model's useful variation while providing a
smaller space to explore. Whether it also improves a cosmological analysis
is a separate question that requires testing against lensing observables.

We investigate the representation question using synthetic alignment-response
surfaces. Each surface gives the response on a grid of wavenumber $k$,
which labels spatial scale, and redshift $z$, which tracks cosmic expansion.
We extend a nonlinear alignment (NLA) prescription with smooth factors controlling
redshift and scale dependence \citep{2007NJPh....9..444B}. Thirteen shape
parameters define the family; the overall normalization and a prescribed
cosmological factor remain separate. We check the numerical behaviour of
the model across its parameter ranges and generate 100,000 surfaces, each
containing 3,131 grid values rather than 3,131 independent physical parameters.

We compare three ways to encode each surface with a short vector of
\emph{latent coordinates}. Principal component analysis (PCA) reconstructs
surfaces as linear combinations of shared patterns
\citep{2016RSPTA.37450202J}. A convolutional autoencoder learns nonlinear
encoding and reconstruction directly on the grid, while a PCA-assisted
autoencoder first projects the surface onto a fixed linear basis and then
compresses its coefficients \citep{2006Sci...313..504H}. The comparison uses
the same 80,000 training and 10,000 validation surfaces at latent dimensions
$L\in\{2,4,6,8,10\}$. Preprocessing is fitted on training data. We examine
average reconstruction errors, difficult surfaces, and associations between
the learned coordinates and the generating parameters. The remaining
10,000 samples are reserved for final test evaluation after model choices
are fixed.

PCA and autoencoders are established methods, and neural networks already
accelerate cosmological calculations \citep{2022MNRAS.511.1771S}. The
contribution here is their controlled application to a deliberately flexible
IA response family: interpretable model construction is linked to a
comparison at equal representation length and to diagnostics of the
approximation's limits. Six Conv1D coordinates reach our reference target
of $99.9\%$ recovered validation variance, with substantially lower mean
relative error than six-component PCA. The remaining local errors show why
variance recovery alone is insufficient. This combination of a compact
representation and explicit failure diagnostics provides a basis for
further tests with observables.

Section~\ref{sec:background} introduces the cosmological background and
explains how intrinsic alignment enters galaxy-shape correlations.
Section~\ref{sec:model} defines the response model and its sampling.
Section~\ref{sec:autoencoder} describes the compression methods, training,
and accuracy measures. Section~\ref{sec:results} compares the results and
examines their limitations, and Section~\ref{sec:summary} draws the
methodological conclusions and sets out the steps towards survey use.
